\expandafter\def\csname HistoricalNoteText1\endcsname{Sub quantitate imaginaria hic semper intelligo quantitatem in forma $a+b\sqrt{-1}$ contentam, quamdiu $b$ non est $=0$. In hoc sensu expressio illa semper ab omnibus geometris primae notae accepta est, neque audiendos censeo, qui quantitatem $a+b\sqrt{-1}$ in eo solo casu imaginaria vocare voluerunt ubi $a=0$,
impossibilem vero quando non sit $a=0$, quum haec distinctio neque necessaria sit neque ullius utilitatis. ---
Si quantitates imaginariae omnino in analysi retineri debent (quod pluribus rationibus consultius videtur, quam
ipsas abolere, modo satis solide stabiliantur): necessario tamquam aeque possibiles ac reales spectandae sunt;
quamobrem reales et imaginarias sub denominatione communi \textit{quantitatum possibilium} complecti mallem: contra, \textit{impossibilem} dicerem quantitatem, quae conditionibus satisfacere debeat, quibus ne imaginariis quidem
concessis satisfieri potest, attamen ita, ut \textit{phrasis} haec idem significet ac si dicas, talem quantitatem in toto
magnitudinum ambitu non dari. Hinc vero genus peculiare quantitatum formare, neutiquam concederem.
Quodsi quis dicat, triangulum rectilineum aequilaterum rectangulum impossibile esse, nemo erit qui neget.
At si tale triangulum impossibile tamquam novum triangulorum genus contemplari, aliasque triangulorum proprietates ad illud applicare voluerit, ecquis risum teneat? Hoc esset verbis ludere seu potius abuti. --- Quamvis vero etiam summi mathematici saepius veritates, quae quantitatum ad quas spectant possibilitatem manifeste supponunt, ad tales quoque applicaverint quarum possibilitas adhuc dubia erat; neque abnuerim, huiusmodi licentias plerumque ad solam formam et quasi velamen ratiociniorum pertinere, quod veri geometrae
acies mox penetrare possit: tamen consultius, scientiaeque, quae tamquam perfectissimum claritatis et certitudinis exemplar merito celebratur, sublimitate magis dignum videtur, tales libertates aut omnino proscribere,
aut saltem parcius neque alias ipsis uti, nisi ubi etiam minus exercitati perspicere valeant, rem etiam absque
illarum subsidio etsi forsan minus breviter tamen aeque rigorose absolvi potuisse. --- Ceterum haud negaverim,
ea quae hic contra impossibilium abusum dixi, quodam respectu etiam contra imaginarias obiici posse: sed harum vindicationem nec non totius huius rei expositionem uberiorem ad aliam occasionem mihi reservo.}

\expandafter\def\csname HistoricalNoteText2\endcsname{Observare convenit, ill. \textsc{d'Alembert} in sua huius demonstrationis expositione considerationes geometricas adhibuisse, atque $X$ tamquam abscissam, $x$ tamquam ordinatam curvae spectavisse (secundum morem
omnium geometrarum primae huius saeculi partis, apud quos notio functionum minus usitata erat). Quia vero
omnia ipsius ratiocinia, si ad ipsorum essentiam solam respicis, nullis principiis geometricis, sed pure analyticis innituntur, et curvae imaginariae, ordinataeque imaginariae expressiones duriores esse lectoremque hodiernum facilius offendere posse videntur, formam repraesentationis mere analyticam hic adhibere malui. Hanc
annotationem ideo adieci, ne quis demonstrationem \textsc{d'Alembertianam} ipsam cum hac succincta expositione comparans aliquid essentiale immutatum esse suspicetur.}

\expandafter\def\csname HistoricalNoteText3\endcsname{Hacce occasione obiter adnoto, ex harum serierum numero plurimas esse, quae primo aspectu maxime
convergentes videantur, e. g. ad maximam partem eas, quibus ill. \textsc{Euler} in parte poster. \textit{Inst. Calc. Diff.}
Cap. VI. ad summam aliarum serierum quam proxime assignandam utitur p. 441--474 (reliquae enim series
p. 475--478 revera convergere possunt), quod, quantum scio, a nemine hucusque observatum est. Quocirca
magnopere optandum esset, ut dilucide et rigorose ostenderetur, cur huiusmodi series, quae primo citissime,
dein paullatim lentius lentiusque convergunt, tandemque magis magisque divergunt, nihilominus summam
proxime veram suppeditent, si modo non nimis multi termini capiantur, et quousque talis summa pro exacta
tuto haberi possit?}

\expandafter\def\csname HistoricalNoteText4\endcsname{Tota haec res multum illustrabitur per aliam disquisitionem sub prelo iam sudantem, ubi in argumento
longe quidem diverso, nihilominus tamen analogo, licentiam similem prorsus eodem iure usurpare
potuissem, ut hic in aequationibus ab omnibus analystis factum est. Quamquam vero plurium veritatum
demonstrationes adiumento talium fictionum paucis verbis absolvere licuisset, quae absque his perquam difficiles
evaserunt et subtilissima artificia requirunt, tamen illis omnino abstinere malui, speroque, paucis me satisfacturum
fuisse, si analystarum methodum imitatus essem.}

\expandafter\def\csname HistoricalNoteText5\endcsname{E. per errorem habet $C$, unde etiam postea perperam statuit $pqr=C$.}

\expandafter\def\csname HistoricalNoteText6\endcsname{In hanc expositionem error irrepsisse videtur, scilicet p. 118. l. 5. loco characteris $p$ (on choisissait
seulement celles où entrait $p$ etc.), necessario legere oportet, \textit{une même racine quelconque de l'équation
proposée}, aut simile quid, quum illud nullum sensum habeat.}
